% Report-ready section. Requires amsmath, amsthm, and booktabs.
% The main report should cite the pinned PA194 and PA195S/C archives.
\section{The cached-head construction and its correctness}
\label{sec:new-cached-proof}

The new 194-state machine combines the supplied PA194 axiom decoder
with the cached-head command protocol previously realized at 195 states.
The logical calculus, the 190-token separate-head axiom emitter, the
physical register layout, and the arithmetic engine are retained.
The contribution is the verified composition at 194 states. In particular,
we do not claim that the cached-head protocol itself is new.

In native labels, the six controller changes are
\[
\begin{array}{c|c@{\qquad}c|c}
(q,b)&\text{new action}&(q,b)&\text{new action}\\\hline
(38,0)&0R11&(38,1)&1R122\\
(164,1)&1R130&(214,0)&0R312\\
(214,1)&1R20&(312,0)&0R188.
\end{array}
\]
They are the exact six changes in the supplied 195C builder; they are not
portable to an unrelated machine merely because it uses some of the same
integer labels. The source replacements are the corresponding two
non-axiom subtrees. Applying those changes to the supplied 195S grammar
reconstructs the literal supplied 195C native table, whose SHA-256 begins
\texttt{8f34d9e5cc2e64ecc49bbc8581dc29b2}. The complete hash is in the
machine-readable provenance.

\subsection{The exact command interface}

Let $p(a,b)=(a+b)(a+b+1)/2+b$, write $[a;w]=p(a,w)$, and use total
Cantor projections with $\operatorname{hd}(0)=\operatorname{tl}(0)=0$.
Thus $[0;0]=0$; the proof never assumes that pushing zero increases a
stored list length. Let
\[
 I(a,b)=4p(a,b),\qquad U(a,B)=4p(B,a)+2,
 \qquad\Gamma(B)=I(B,U(B,B)).
\]
For a store $\langle t,w,d,h\rangle$, the cached-head Shift is
\[
 S\langle t,w,d,h\rangle=
 \begin{cases}
 \mathsf{halt},&h=1,\\
 \langle h,\operatorname{tl}(w),d,\operatorname{hd}(w)\rangle,&h\ne1.
 \end{cases}
\]
Putting $y=d+x$, its MP command is
\[
 \mathsf{MP}_x\langle t,w,d,h\rangle=
 \begin{cases}
 \langle0,w,0,y\rangle,&I(t,y)=h,\\
 \langle0,w,y,0\rangle,&I(t,y)\ne h.
 \end{cases}
\]
Neither MP branch pushes. The old 194-state protocol instead tests
$t+h=1$ and returns a successful MP result in the top while pushing
zero. The two source protocols therefore differ operationally. For
example, Shift with $t=1,h=0$ halts in the old protocol and continues
in the cached-head protocol. They implement the same refutation search
by different proof-command compilers.

The separate-head AX command pushes ten entries including the saved top,
clears $d,h$, and leaves the guarded old gate as the new top. Its nine
generated entries, from top downward, are
\[
 \Gamma(c),L_6,L_5,L_4,L_3,L_2,L_1,
 \rho_8(d+x,c,h+y,d),U(d,h).
\]
The gate is zero when its numerical guard fails. Every generated entry
other than the last has residue zero modulo four; the last has residue
two. The gate also has residue zero. None is code one, whether or not the
guard succeeds. This syntactic fact is sufficient for the termination
assertions below.

\subsection{A uniform compiler for the changed protocol}

It is useful to state the stack identities for a general AX push count
$n$, then specialize to $n=10$. Define
\begin{align*}
 \mathsf{Push}_n&=\mathsf{AX}_{0,0,0};S^{n-1};\mathsf{MP}_0,\\
 \mathsf{Pop}_n&=\mathsf{AX}_{0,0,0};S^{n+1},\\
 \mathsf{Discard}_n&=\mathsf{MP}_1;\mathsf{AX}_{1,0,0};S^n;\mathsf{MP}_0.
\end{align*}
The total pairing projection identities give
\begin{align*}
 \mathsf{Push}_n\langle t,w,d,h\rangle
  &=\langle0,[t;w],0,0\rangle,\\
 \mathsf{Pop}_n\langle t,[a;w],d,h\rangle
  &=\begin{cases}
    \mathsf{halt},&t=1,\\
    \langle t,w,0,a\rangle,&t\ne1,
    \end{cases}\\
 \mathsf{Discard}_n\langle t,w,0,0\rangle
  &=\langle0,w,0,0\rangle.
\end{align*}
For Push, the final MP comparison fails because an implication code
cannot equal the generated universal code. Pop tests the preserved top
only at its last Shift. For Discard, $I(t,1)>0$ forces the first MP to
fail, and $I(U(1,0),0)>0$ forces the last to fail. These arguments also
cover the zero-stack alias. Pop is deliberately not total when the
preserved top is one.

Let $\mathsf{Offset}_0$ be empty and
$\mathsf{Offset}_i=\mathsf{MP}_i$ for $i>0$. Generalization is compiled by
\[
 \mathsf{Gen}_{i,n}=
 \mathsf{Push}_n;\mathsf{Offset}_i;S;\mathsf{AX}_{i,0,0};S^n.
\]
It takes $\langle a,w,0,0\rangle$ to
$\langle U(i,a),w,0,0\rangle$ for every $a,w$, including $a=1$.
For an ordinary generated head in top-down position $j$, with
$1\le j\le n-3$, use
\[
 \mathsf{Discard}_n;\mathsf{AX}_{A,D,H};S^{j+1};
 \mathsf{Pop}_n^{n-1-j}.
\]
It returns that head with the original underlying stack and zero offset
and cache. The separate diagonal head has $j=1$. The restricted
propositional head with $u\le a$ is selected by
\[
 \mathsf{Discard}_n;\mathsf{Offset}_u;
 \mathsf{AX}_{b,a-u,c};S^{n-1};\mathsf{Pop}_n,
\]
and the gate at $(a,b,c)$ by
$\mathsf{Discard}_n;\mathsf{AX}_{c,b,a};\mathsf{Pop}_n^n$.
All Pop uses in these selectors preserve a head different from one.

The key changed MP suffix is
\[
 \mathsf{MPSuffix}_b=S;S;\mathsf{MP}_b;S.
\]
Starting from $\langle0,[a,I(a,b),0;w],0,0\rangle$, this word halts
if $a=1$ or $b=1$; otherwise it returns
$\langle b,w,0,0\rangle$. The dummy zero in its precondition is
essential: the final Shift consumes it, so the arbitrary underlying
stack is preserved. The same identity holds when all three entries
alias the zero stack.

For a proof tree $p$ whose last step is MP, with minor proof $p_0$ of
$a$ and major proof $p_1$ of $I(a,b)$, define
\[
\begin{split}
 W_p={}&\mathsf{Discard}_n;\mathsf{Push}_n;
 W_{p_1};\mathsf{Push}_n;\\
 &W_{p_0};\mathsf{Push}_n;\mathsf{MPSuffix}_b.
\end{split}
\]
Axiom proofs use the selectors above, and generalization uses
$W_{p_0};\mathsf{Gen}_{i,n}$. Induction on $p$ proves that from every
$\langle t,w,0,0\rangle$, the word $W_p$ either halts or returns
$\langle\operatorname{concl}(p),w,0,0\rangle$. An early halt is harmless
for completeness. Therefore, if $p$ proves one,
\[
             \mathsf{AX}_{0,0,0};W_p;\mathsf{Pop}_{10}
\]
halts from every incoming store. The initial AX clears the offset and
cache, and the final Pop tests the proved top before inspecting the
arbitrary underlying stack. This universal incoming-store property is
the correct hypothesis for the inherited exhaustive command scheduler;
it does not assume that earlier trials restore an empty store.

\subsection{Soundness and the changed primitive tails}

Annotate the top, cache, and every canonical live stack entry with
derivations in the numerical calculus. No derivation is required for
the offset $d$. Code zero has a fixed derivation, such as the all-zero
dominated propositional instance. AX preserves the annotation invariant
by its axiom suppliers, generalization, and saved top. MP preserves it
by modus ponens or by the fixed zero proof. Shift transfers existing
proofs; a halting Shift exposes a derivation of code one. This proves
soundness of the cached protocol independently of the original protocol.

The changed source tails have uniform elementary loop proofs. With
$N=3649$, the successful MP tail begins at $N-15=3634$ and transfers
$(D,H)=(y,0)$ to $(0,y)$ by the loop
``decrement $D$; increment $H$; return to the loop''. It exits after
$3y+1$ instructions, then returns to fetch in one more instruction.
Failure returns directly with $(D,H)=(y,0)$.
The new Shift begins at $N-11=3638$. Under the inherited entry condition
$X=0$, it clears $T$, copies $H$ into $T,X$, and restores $H$ from $X$.
It reaches $(T,H,X)=(h,h,0)$ in $2t+7h+3$ instructions. Two more
instructions set $X=1$ and enter the unchanged guard and Cantor decoder.
Thus the primitive program realizes the stated cache-only test.

The source uses the same AX emitter, numerical calculus, and scheduler
as the supplied cached parent. The resolved source maps of the new
decoder and that exact parent agree at all 3,649 primitive sites,
including every successor. Their finite source traces therefore agree
after identifying counters by instruction ordinal. This equality is a
fixed finite fact and its trace consequence is provable in PA.

\subsection{Why the physical decoder transports to the cached source}

The raw allocation is 134 fixed controls plus 67 AX readers. The seven
PA194 folds remain compatible, giving
\[
                         134+67-7=194.
\]
Their local detours inspect only the two relevant tape cells and use
none of the six modified controller rows. The new controller also adds
no incoming edge to an exceptional fold control. More directly, all
28 history/read-bit cases are rechecked in the final simultaneous table,
and all 23 incoming-context rows establish the required history bit.
This is a quotient of the newly compiled cached raw machine; the proof
does not pretend that the new source graph equals the old uncached one.

Both symbolic implementations check all 4,749 source outcomes for
arbitrary natural register values, both decrement branches, an arbitrary
nonempty finite spare list, and the affine gap parameter. The independent
implementation rederives its four generic list-operation lemmas from
the actual rows and uses no arithmetic-box oracle. It records a minimum
positive-duration constant of three. The whole blank tape reaches the
exact clean source frame after 57,922 transitions in both implementations.

\begin{theorem}[PA-internal preservation]
Let $H_C$ and $H_{195C}$ be the finite blank-tape halting sentences of
the new 194-state table and the supplied cached-head 195-state table.
Then
\[
                 \mathrm{PA}\vdash H_C\leftrightarrow H_{195C}.
\]
\end{theorem}
\begin{proof}
Let $R$ be their common finite register program with the fixed clean
seed. Each nonhalting source instruction has an exact positive-length
physical simulation to the next frame, and a source halt reaches the
terminal control. The certificate operations are literal rows, finite
concatenations, constant-bit scans, and unary-list register operations.
Their concrete witness traces are primitive recursively constructible;
PA proves the scan and update lemmas by induction on finite lengths.
Replaying the finite certificate list gives those statements uniformly
in coded frames, rather than only for sampled standard inputs.

A finite halting source trace produces a physical trace by prepending
the verified startup and concatenating its segments. Conversely, suppose
a physical run first halts after $T$ steps. If the first $T+1$ source
instructions did not halt, their positive segments, together with
startup, would give a physical run longer than $T$ ending in a nonhalt
frame. Determinism contradicts the alleged first halt. Thus the bounded
source iteration contains a halt. These finite recursions and the
determinism comparison are provable in PA. Applying the argument to
both tables and eliminating the common source-halting predicate proves
the claim.
\end{proof}

Composing this preservation theorem with the inherited PA calibration
of the separate-head calculus and the cached-command compiler gives
\[
 \boxed{\mathrm{PA}\vdash
   \bigl(\neg H_C\leftrightarrow\operatorname{Con}(\mathrm{PA})\bigr).}
\]
The new result is an ordinary mathematical proof supported by executable
machine-specific certificates. No newly compiled end-to-end Lean or Rocq
theorem for this literal table is asserted.
